% Generated draft: requires mathematical and source review.
\documentclass[11pt]{article}
\usepackage[margin=25mm]{geometry}
\usepackage{fontspec}
\setmainfont{Latin Modern Roman}
\usepackage{amsmath,amssymb,mathtools,unicode-math}
\setmathfont{Latin Modern Math}
\usepackage{enumitem,needspace,xurl,hyperref}
\setlength{\parindent}{0pt}
\setlength{\parskip}{5pt}
\setlength{\emergencystretch}{3em}
\newcommand{\N}{\mathbb{N}}
\newcommand{\Z}{\mathbb{Z}}
\newcommand{\Q}{\mathbb{Q}}
\newcommand{\R}{\mathbb{R}}
\newcommand{\C}{\mathbb{C}}
\newcommand{\F}{\mathbb{F}}
\newcommand{\A}{\mathbb{A}}
\newcommand{\PP}{\mathbb{P}}
\newcommand{\E}{\mathbb{E}}
\newcommand{\eps}{\varepsilon}
\newcommand{\dd}{\,\mathrm{d}}
\begin{document}
\section*{UnsolvedMath ID: 2100605}
\textbf{Source problem code:} catalogue entry
\subsection{Definitions and mathematical statement}
Let $(M^{2n},\pi)$ be a Poisson manifold: $\pi$ is a bivector field whose induced bracket $\{f,g\}=\pi(df,dg)$ satisfies the Jacobi identity. Assume its nondegenerate locus
\[
 U=\{x\in M:\operatorname{rank}\pi_x=2n\}
\]
is dense; this locus is open and carries the inverse symplectic form. No compactness or connectedness hypothesis on $M$ is supplied.

Under the usual toric convention, the action is an effective action of an $n$-torus $T^n$, with Lie algebra $\mathfrak t$, preserving the Poisson structure. Hamiltonianity on $U$ means that there is a moment map $\mu_U:U\to\mathfrak t^*$ such that, with a fixed Hamiltonian sign convention,
\[
 \xi_M|_U=\pi^{\sharp}\bigl(d\langle\mu_U,\xi\rangle\bigr)
 \quad\text{for every }\xi\in\mathfrak t.
\]
The record's word ``toric'' is not further defined, so effectiveness and the precise torus dimension should be confirmed against the full source.

\textbf{Question.} For such an even-dimensional Poisson manifold with dense symplectic locus and a toric action Hamiltonian on that locus, is there an analog of Delzant's theorem for the image of the moment map?

In the classical compact connected symplectic toric setting, Delzant's theorem relates the system to its moment polytope. Here the literal assumptions guarantee a moment map only on $U$, so its presently specified image is $\mu_U(U)\subset\mathfrak t^*$. They do not supply an extension $\mu:M\to\mathfrak t^*$, properness, a type of Poisson degeneracy, or a class of generalized polytopes. It is also unclear whether the requested analog is a theorem describing possible images, a convexity theorem, or a classification/reconstruction theorem for the toric Poisson manifolds themselves. No such extra target or hypothesis is silently selected.

The supplied research summary reports that it could not verify a published Delzant-type classification in this general dense-symplectic Poisson setting during that session, and notes possible degeneracies or singular features of moment images. This is an open/unverified evidence assessment, not an independently established claim that no such theorem exists.
\subsection{Short English statement}
For a generically symplectic Poisson manifold with a torus action Hamiltonian on its symplectic locus, seek a Delzant-type description of the moment-map image. The source leaves global extension of the moment map, compactness, degeneracy type, and the desired classification statement unspecified.
\subsection{Sources}
\begin{itemize}
\item Local UnsolvedMath registry, permanent ID 2100605, source code catalogue entry, ``Open Problems in Integrable Systems --- Poisson geometry and action-angle variables''; supplied status open, classification OPEN-TRIAGE, updated 2026-08-06.
\item \emph{Open Problems, Questions, and Challenges in Finite-Dimensional Integrable Systems} (2018), Problem 6.5, \url{https://arxiv.org/abs/1804.03737}. The record reports source-TeX extraction and access on 2026-07-29. It says the source presents an open item and current status requires release review. Rights: NEEDS\_REVIEW; public arXiv source TeX inspected according to the record.
\item Delzant's theorem is named in the supplied statement and summary without a bibliographic reference. The summary gives no verified citation completing the proposed Poisson analog.
\end{itemize}
\end{document}
